% Modelo entidad-relación de PayCore (simplificado a lo esencial).
\begin{tikzpicture}[
  font=\scriptsize,
  >={Stealth[length=2mm]},
  entidad/.style={draw=acento!55, fill=acento!7, rounded corners=3pt,
                  text width=40mm, inner sep=5pt},
  titulo/.style={font=\scriptsize\bfseries, color=acento},
  rel/.style={->, draw=tinta!55},
  nota/.style={font=\scriptsize\itshape, color=tinta!65},
]

\node[entidad] (users) {%
  \textcolor{acento}{\bfseries users}\\[2pt]
  \codigo{id} \textsc{uuid pk}\\
  \codigo{email} \textsc{citext uk}\\
  \codigo{password\_hash}\\
  \codigo{full\_name}\\
  \codigo{created\_at}};

\node[entidad, right=26mm of users] (accounts) {%
  \textcolor{acento}{\bfseries accounts}\\[2pt]
  \codigo{id} \textsc{uuid pk}\\
  \codigo{user\_id} \textsc{fk uk}\\
  \codigo{alias} \textsc{uk}\\
  \codigo{kind} user\,|\,system\\
  \codigo{currency} = PEN\\
  \codigo{balance\_cents} \textsc{bigint}\\
  \codigo{created\_at}};

\node[entidad, below=14mm of users] (idem) {%
  \textcolor{acento}{\bfseries idempotency\_keys}\\[2pt]
  \codigo{id} \textsc{uuid pk}\\
  \codigo{user\_id} \textsc{fk}\\
  \codigo{key}\\
  \codigo{request\_fingerprint}\\
  \codigo{state}\\
  \codigo{response\_status}\\
  \codigo{response\_body} \textsc{jsonb}\\
  \codigo{transfer\_id} \textsc{fk}\\
  \codigo{expires\_at}};

\node[entidad, below=14mm of accounts] (transfers) {%
  \textcolor{acento}{\bfseries transfers}\\[2pt]
  \codigo{id} \textsc{uuid pk}\\
  \codigo{source\_account\_id} \textsc{fk}\\
  \codigo{destination\_account\_id} \textsc{fk}\\
  \codigo{amount\_cents} \textsc{bigint}\\
  \codigo{status}\\
  \codigo{rejection\_reason}\\
  \codigo{concept}\\
  \codigo{completed\_at}};

\node[entidad, right=22mm of transfers] (ledger) {%
  \textcolor{acento}{\bfseries ledger\_entries}\\[2pt]
  \codigo{id} \textsc{uuid pk}\\
  \codigo{transfer\_id} \textsc{fk}\\
  \codigo{account\_id} \textsc{fk}\\
  \codigo{amount\_cents} \textsc{bigint}\\
  \hspace{4pt}{\itshape signo: $-$ débito, $+$ crédito}\\
  \codigo{balance\_after\_cents}\\
  \codigo{created\_at}};

\draw[rel] (users.east) -- node[above] {1\,:\,1} (accounts.west);
\draw[rel] (users.south) -- node[right] {1\,:\,N} (idem.north);
\draw[rel] (accounts.south) -- node[right] {1\,:\,N} (transfers.north);
\draw[rel] (transfers.east) -- node[above, align=center]
  {1\,:\,2\\[-2pt]{\itshape exactamente}} (ledger.west);
\draw[rel] (idem.east) -- node[above] {0\,:\,1} (transfers.west);

\node[nota, below=5mm of transfers.south, text width=140mm,
      align=left, anchor=north] {%
  El libro mayor (\codigo{ledger\_entries}) es de solo inserción y es la
  fuente de verdad. La columna \codigo{accounts.balance\_cents} es una
  proyección materializada que se escribe en la misma transacción. La
  restricción de no negatividad exime a las cuentas de sistema, que son el
  origen auditado del saldo ficticio.};

\end{tikzpicture}
